\uuid{DrSs}
\exo7id{4436}
\titre{exo7 4436}
\auteur{quercia}
\organisation{exo7}
\datecreate{2010-03-14}
\isIndication{false}
\isCorrection{true}
\chapitre{Série numérique}
\sousChapitre{Autre}
\cpge{mpsi-mp2i.procedes-sommatoires-discrets.a, pcsi.series-numeriques.a, ptsi.series-numeriques.a, mpsi-mp2i.procedes-sommatoires-discrets.b, pcsi.series-numeriques.b, ptsi.series-numeriques.b, mpsi-mp2i.calcul-algebrique-trigonometrie.d, pcsi.calcul-algebrique-trigonometrie.d, ptsi.calcul-algebrique-trigonometrie.d, bcpst1.trigonometrie, bcpst1.suites-reelles-usuelles, mp-mpi.series-numeriques-vectorielles.a, pt.series-numeriques.b, mp-mpi.series-numeriques-vectorielles.b, pc.series-numeriques, psi.series-numeriques.a, pt.series-numeriques.a, bcpst2.series-reelles}
\module{Analyse}
\niveau{L1}
\difficulte{}

\contenu{
\question{
Montrer que
$\sum_{k=0}^\infty \Arctan\Bigl(\frac1{k^2+k+1}\Bigr) = \frac \pi2$.
(On pourra calculer $\tan s_n$)
}
\reponse{
$\tan s_n = n+1$ par récurrence et
         $s_n \le \sum_{k=0}^\infty \frac 1{k^2+k+1} \le 1
         + \sum_{k=0}^\infty \frac1{n(n+1)} = 2$.
}
}
